\begin{table}[H]
\captionsetup{justification=raggedright,singlelinecheck=false}
\centering
\caption{\textit{Gap Analysis}}
\label{tbl:gap-analysis}
\begin{tabular}{|p{2.7cm}|p{3.5cm}|p{3.5cm}|p{3.5cm}|}
\hline
\textbf{Domain Masalah} &
\textbf{Kondisi Saat Ini (\textit{As-Is})} &
\textbf{Kesenjangan (\textit{Gap})} &
\textbf{Kondisi Ideal (\textit{To-Be})} \\
\hline

Aktivitas \textit{crawling} &
Proses \textit{crawling} dijalankan dari \textit{executable} di disk dan menggunakan akses file/direktori untuk mengambil data, sehingga meninggalkan artefak. &
Belum ada mekanisme untuk meminimalkan jejak \textit{crawling}. &
\textit{Crawling} dilakukan dengan jejak minimal melalui \textit{in-memory processing}, penghapusan \textit{buffer} otomatis per \textit{batch}, dan tanpa meninggalkan artefak yang mudah terdeteksi di disk. \\
\hline

Kanal eksfiltrasi dan karakteristik trafik &
Eksfiltrasi menggunakan protokol umum tetapi pola trafik tidak disamarkan dan tetap mudah dibedakan dari komunikasi aplikasi sah. &
Kanal komunikasi belum menyerupai trafik normal dan karakteristik paket belum tersamarkan sehingga aktivitas dapat terdeteksi. &
Eksfiltrasi berlangsung melalui kanal yang menyerupai trafik sah dengan pola komunikasi adaptif, terenkripsi, dan tidak mencolok. \\
\hline

Mekanisme eksfiltrasi &
Data hasil \textit{crawling} langsung dikirim ke \textit{endpoint} tanpa enkripsi atau kompresi khusus, sehingga struktur \textit{payload} masih mudah diinterpretasi. &
Tidak ada mekanisme untuk mengamankan \textit{payload} eksfiltrasi. &
Eksfiltrasi dilakukan dengan \textit{payload} terenkripsi, dikemas, atau disamarkan sehingga tidak mudah dibaca maupun dianalisis. \\
\hline

\end{tabular}
\end{table}
